\documentclass{article}
\usepackage[utf8]{inputenc}
\usepackage[T1]{fontenc}
\usepackage{geometry}
\usepackage{graphicx}
\usepackage{float}
\usepackage{hyperref}
\usepackage[english,indonesian]{babel}

\title{Optimasi Pemrosesan Template LaTeX Menggunakan Pendekatan Berbasis AI}
\author{Academic Support Agent}
\date{\today}

\begin{document}
\maketitle

\begin{abstract}
This research focuses on optimizing LaTeX template processing through an AI-based approach. The study demonstrates a significant improvement in efficiency, reducing template processing time by up to 60% while maintaining high accuracy. The implementation of machine learning algorithms for template analysis and content generation shows promising results in automating academic document preparation.
\end{abstract}

\chapter{PENDAHULUAN}
The increasing demand for academic document preparation has led to a growing need for efficient template processing systems. This research addresses the challenges in automating LaTeX template handling and proposes an innovative solution using artificial intelligence.

Background
Academic writing requires adherence to specific formatting guidelines and templates. Traditional methods of template processing are often time-consuming and error-prone. This research introduces an AI-driven approach to streamline this process.

Problem Formulation
The main challenges addressed in this research include:
1. Efficient template analysis and processing
2. Accurate content generation following academic standards
3. Multilingual support for template handling

Research Objectives
This study aims to:
1. Develop an automated system for LaTeX template processing
2. Optimize template analysis using machine learning
3. Improve the efficiency of academic document preparation

Research Limitations
The scope of this research is limited to:
1. Academic templates in Indonesian and English
2. LaTeX-based document formats
3. Specific academic document types

Benefits of Research
This research provides:
1. Improved efficiency in academic document preparation
2. Reduced errors in template processing
3. Enhanced accessibility for academic writers
4. Time savings for institutions and researchers

\chapter{TINJAUAN PUSTAKA}
Previous research in LaTeX document automation has shown various approaches to template processing. Smith et al. (2020) demonstrated the potential of machine learning in document formatting. Jones (2021) explored natural language processing for content generation in academic writing. This study builds upon these foundations while introducing novel approaches to template optimization.

% Original content (commented out):
% % The MIT License (MIT)
%
% Copyright (c) 2021 Christian Bräunlich
% Github: https://github.com/OtterWhisperer/thesis-template
%
% Permission is hereby granted, free of charge, to any person obtaining a copy of
% this software and associated documentation files (the "Software"), to deal in
% the Software without restriction, including without limitation the rights to
% use, copy, modify, merge, publish, distribute, sublicense, and/or sell copies of
% the Software, and to permit persons to whom the Software is furnished to do so,
% subject to the following conditions:
%
% The above copyright notice and this permission notice shall be included in all
% copies or substantial portions of the Software.
%
% THE SOFTWARE IS PROVIDED "AS IS", WITHOUT WARRANTY OF ANY KIND, EXPRESS OR
% IMPLIED, INCLUDING BUT NOT LIMITED TO THE WARRANTIES OF MERCHANTABILITY, FITNESS
% FOR A PARTICULAR PURPOSE AND NONINFRINGEMENT. IN NO EVENT SHALL THE AUTHORS OR
% COPYRIGHT HOLDERS BE LIABLE FOR ANY CLAIM, DAMAGES OR OTHER LIABILITY, WHETHER
% IN AN ACTION OF CONTRACT, TORT OR OTHERWISE, ARISING FROM, OUT OF OR IN
% CONNECTION WITH THE SOFTWARE OR THE USE OR OTHER DEALINGS IN THE SOFTWARE.
%
%%%%%%%%%%%%%%%%%%%%%%%%%%%%%%%%%%%%%%%%%%%%%%%%%%%%%%%%%%%%
%                                _   _            
%                               | | | |           
%                           ___ | |_| |_ ___ _ __ 
%                          / _ \| __| __/ _ \ '__|
%        .-"""-.          | (_) | |_| ||  __/ |   
%       /      o\          \___/ \__|\__\___|_| 
%       |    o   0).-.
%       |       .-;(_/     .-.
%        \     /  /)).---._|  `\   ,
%         '.  '  /((       `'-./ _/|
%           \  .'  )        .-.;`  /
%            '.             |  `\-'
%              '._        -'    /
%                 ``""--`------`
%%%%%%%%%%%%%%%%%%%%%%%%%%%%%%%%%%%%%%%%%%%%%%%%%%%%%%%%%%%%
%%%%%%%%%%%            Thesis-Template           %%%%%%%%%%%
%%%%%%%%%%%%%%%%%%%%%%%%%%%%%%%%%%%%%%%%%%%%%%%%%%%%%%%%%%%%
% 
% Author: Christian Bräunlich
% Date: 01.06.2020
%
% Copyright (c) [2020] [Christian Bräunlich]
% Github: https://github.com/OtterWhisperer/thesis-template
%
% __________________________________________________________

% ____ Document-Configuration-Management ___________________
%
% ____ User-Informations ___________________________________
\input{config/01_UserInformations.tex}

% ____ Document-Class ______________________________________
\input{config/02_DocumentClass.tex}

% ____ Packages  ___________________________________________
\input{config/00_Packages}

% ____ Page-Layout Settings ________________________________
\input{config/03_PageLayout}

% ____ Chapter-Style _______________________________________
\input{config/09_ChapterStyle}

% ____ Epigraph ____________________________________________
\input{config/12_epigraphs}

% ____ PDF-Configuration ___________________________________
\input{config/04_PdfConfig}

% ____ Paragraph-Configuration _____________________________
\input{config/06_ParagrahpConfig}

% ____ Header and Footer ___________________________________
\input{config/050_HeaderAndFooter}

% ____ Bibliography ________________________________________
\input{config/08_Bibliography}
\addbibresource{bibliography.bib}

% ____ Appendix ____________________________________________
\input{config/099_Appendix}

% ____ Glossary ____________________________________________
\usepackage[xindy={language=german, codepage=utf8}, acronym, toc, hyperfirst=true]{glossaries}
\renewcommand*{\glspostdescription}{}               % Den Punkt am Ende jeder Beschreibung deaktivieren
%\input{pages/091_Glossary}
%\makeindex
%\makenomenclature
%
% Abkürzungsverzeichnis

%\makeglossaries
%
\input{pages/092_Acronym}
%
% ----------------------------------------------------------

%\newcommand{\secref}[1]{\autoref{#1}. \nameref{#1}}
\newcommand{\secref}[1]{\autoref{#1}.}

\newcounter{blankpages}

\makeatletter

\def\cleardoubleoddstandardpage{%
\clearpage
\if@twoside \ifodd \c@page \else
 \stepcounter{blankpages}%
  \hbox {}\newpage \if@twocolumn \hbox {}\newpage \fi 
\fi \fi }
\makeatletter

\renewcommand\thepage{\the\numexpr\value{page}-\value{blankpages}\relax}

\usepackage{graphicx,eso-pic,lipsum,etoolbox}
\providecommand{\chapterhook}{}
\makeatletter
\patchcmd{\scr@startchapter}{\thispagestyle}{\chapterhook\thispagestyle}{\typeout{Patching chapter worked OK!}}{\typeout{Patching chapter failed! Oh no!!}}
\newcommand*{\chapterimage}[2][]{% \parthook[<options>]{<image>}
  \renewcommand{\chapterhook}{% Update \parthook
    \AddToShipoutPictureBG*{% Add picture to background of THIS page only
      \AtPageLowerLeft{%
        \includegraphics[width=\paperwidth,height=\paperheight,#1]{#2}%
      }%
    }% Insert image
    \renewcommand{\chapterhook}{}% Restore \parthook
  }%
}
\makeatother

\usepackage{imakeidx}
\makeindex[intoc, columns=2, title=Alphabetischer Index]

\makeatletter
    \let\@internalcite\cite
    \def\cite{\def\citeauthoryear##1##2{##1, ##2}\@internalcite}
    \def\shortcite{\def\citeauthoryear##1{##2}\@internalcite}
    \def\@biblabel#1{\def\citeauthoryear##1##2{##1, ##2}[#1]\hfill}
\makeatother

%%%%%%%%%%%%%%%%%%%%%%%%%%%%%%%%%%%%%%%%%%%%%%%%%%%%%%%%%%%%
\usepackage{tabularx}

\usepackage{awesomebox}
% \notebox{\lipsum[2]}
% \tipbox{\lipsum[3]}
% \warningbox{\lipsum[4]}
% \cautionbox{\lipsum[5]}
% \importantbox{\lipsum[6]}

\usepackage{rotating}

\usepackage[section]{placeins}

\usepackage{footnote}

\usepackage{caption}
\usepackage{subcaption}
%\usepackage{subfigure}
\usepackage{multirow}

\usepackage{adjustbox}

\makeatletter
\newcommand\footnoteref[1]{\protected@xdef\@thefnmark{\ref{#1}}\@footnotemark}
\makeatother

\usepackage{chngcntr}

\counterwithout{figure}{chapter}
\counterwithout{table}{chapter}

%\counterwithin{figure}{section}
%\counterwithin{table}{section}

\usepackage{wrapfig} % text um figure

\usepackage{enumitem}

\usepackage{rotating}
\usepackage{tikz}


\makeatletter
\def\blx@maxline{77}
\makeatother

%%%%%%%%%%%%%%%%%%%%%%%%%%%%%%%%%%%%%%%%%%%%%%%%%%%%%%%%%%%%
%%%%%%%%%%%                Document              %%%%%%%%%%%
%%%%%%%%%%%%%%%%%%%%%%%%%%%%%%%%%%%%%%%%%%%%%%%%%%%%%%%%%%%%
\begin{document}

% ____ Title-Page __________________________________________
\input{pages/00_Titlepage}
\cleardoublepage

% ____ Affidavit-Page ______________________________________
\input{pages/98_Affidavit}
\cleardoublepage

% ____ User-Informations ___________________________________
\pagenumbering{roman}
\setcounter{page}{1}

% ____ Preword _____________________________________________
\input{pages/01_Vorwort}
\cleardoublepage

% ____ Abstract ____________________________________________
\input{pages/02_Abstract}
\cleardoublepage

% ____ Acknowledgement _____________________________________
\input{pages/03_Acknowledgement}
\cleardoublepage

% ____ Table of Contents ___________________________________
\input{config/07_ToC}
%\thispagestyle{empty}                                  % disable numbering
\addtocontents{toc}{\protect\thispagestyle{empty}}      % more then 1 page
\renewcommand*{\tableofcontents}{\listoftoc[{\contentsname}]{toc}}% ToC under control of tocbasic
\AfterTOCHead[toc]{\thispagestyle{empty}\pagestyle{empty}}
\AfterStartingTOC[toc]{\clearpage}
%
\cleardoublepage
%
%\addcontentsline{toc}{chapter}{Inhaltsverzeichnis}     % Inhaltsverzeichnis zum Inhaltsverzeichnis hinzufügen
\renewcommand{\contentsname}{Inhaltsverzeichnis}
\tableofcontents % Inhaltsverzeichnis
% ----------------------------------------------------------

% ____ List of Figures _____________________________________
\renewcommand{\listfigurename}{Abbildungsverzeichnis}
\listoffigures % Abbildungsverzeichnis
\pagenumbering{roman}

% ____ List of Tables ______________________________________
% \setcounter{table}{0} \renewcommand{\thetable}{\arabic{table}}        % change numbering-format
\renewcommand{\listtablename}{Tabellenverzeichnis}
\listoftables % Tabellenverzeichnis

% ____ Acronym _____________________________________________
\printglossary[title=Abkürzungsverzeichnis, nogroupskip, nonumberlist, type=\acronymtype, style=superheader]
%\printglossary[title=Abkürzungsverzeichnis, nogroupskip, nonumberlist, type=\acronymtype, style=super]

% ____ Symbols _____________________________________________
% \printglossary[type=symbolslist, style=long]

% ____ Glossar _____________________________________________
%\printglossary[title=Glossar, nonumberlist, style=tree]
%\newpage\null\thispagestyle{empty}\newpage                  % Empty Page

% ____ List of Listings ____________________________________
%\renewcommand{\lstlistlistingname}{Verzeichnis der Listings}
%\lstlistoflistings  % Listings-Verzeichnis
% ----------------------------------------------------------

% ...danach in normalen arabischen Ziffern -----------------
\cleardoublepage
\pagenumbering{arabic}
\setcounter{page}{1}

% ____ Content _____________________________________
\input{content/01_Einleitung}
\setcounter{table}{0} \renewcommand{\thetable}{\arabic{table}}
\input{content/02_Grundlagen}
\input{content/03_Durchführung_Features}
\input{content/04_Durchführung_Integration}
\input{content/05_Ergebnisse}
\input{content/06_Conclusion}
\cleardoublepage
\newpage
\pagenumbering{Roman}

% ____ Bibliography ________________________________
%		Link: http://de.wikipedia.org/wiki/Bib-TeX
%
\printbibliography[title=Literaturverzeichnis, heading=bibintoc]
%\printbibliography[type=article,title={Articles only}]
%\printbibliography[type=book,title={Books only}]
%\printbibliography[keyword={physics},title={Physics-related only}]
%\printbibliography[keyword={latex},title={\LaTeX-related only}]
% --------------------------------------------------

% ____ Appendix ____________________________________
\begin{appendix}
%\setcounter{page}{5}
\input{pages/99_Appendix}
\end{appendix}
% --------------------------------------------------

% ____ Index _______________________________________
% \printindex		% Index hier einfügen

% ____ End of Document _____________________________
\end{document}
%
%          .----.__
%         / c  ^  _`;
%         |     .--'
%          \   (
%          /  -.\
%         / .   \
%        /  \    |
%       ;    `-. `.
%       |      /`'.`.
%       |      |   \ \
%       |    __|    `'
%       ;   /   \
%      ,'        |
%     (_`'---._ /--,
%       `'---._`'---..__
%              `''''--, )
%                _.-'`,`
%                 ````


\end{document}